%% Appendix E. Where the protocol's counts come from.
%% Moved out of Section VII, where it ran to 1100 words of derivation inside the argument. The
%% body now states what the derivation concludes and points here; nothing was dropped on the way,
%% including the counts an earlier draft quoted and this survey withdrew, which are kept so that a
%% reader can see which test was the wrong one and why.
%% Converted from paper/APPENDIX_E.md for the LaTeX edition. The markdown edition in paper/ is a
%% separate edition of the same prose and is not written by this file.
\section{Where the protocol's counts come from}
\label{app:protocol}

Every count in Table~\ref{tab:protocol} is derived below, and each axis is derived for the
statistic that axis actually reports rather than by one convention applied to all of them: a single
rate takes a Wilson half-width, a matched comparison takes McNemar, a ratio takes the standard
error of the log ratio, and a correlation takes the Fisher-z interval.
Section~\ref{subsec:protocol} states what these derivations conclude. This survey re-ran no method,
so every count here rests on an interval width, a power calculation or another paper's measurement,
and never on a measurement of our own.

Fix the width first, then read off the count. Take a 95 percent Wilson interval on a single
reported rate, at the worst case of $\hat{p} = 0.5$. Writing the half-width of \eqref{eq:wilson} as
\begin{equation}
w(n,\hat{p}) = \frac{z\sqrt{4n\hat{p}\,(1-\hat{p}) + z^{2}}}{2\,(n + z^{2})},
\label{eq:halfwidth}
\end{equation}
the count a target width $w_{0}$ demands is
$n_{\min}(w_{0}) = \min\{\,n : w(n,0.5) \le w_{0}\,\}$, which gives
\begin{equation*}
\begin{array}{lrrrr}
\toprule
\text{half-width } w_{0} \text{ (points)} & 20 & 15 & 10 & 5\\
\midrule
n_{\min} \text{ (trials)} & 21 & 39 & 93 & 381\\
\bottomrule
\end{array}
\end{equation*}
Ten points is the coarsest width at which a single rate is worth printing, so the absolute-rate
minimum is 93, rounded to 100. At 100 trials a reported 80 percent has an interval of 71 to 87
percent, and a reported 50 percent has 40 to 60. A comparison is a different question and a harder
one: two rates each carrying $\pm 10$ points do not resolve a 10-point difference between them,
because the difference's standard error is larger by a factor of $\sqrt{2}$, so the width argument
sets a floor on what is worth reporting and not on what can be compared.

For the A/B comparison the relevant calculation is power, and the design is paired.
Table~\ref{tab:protocol} matches initial conditions by image overlay and interleaves the two
policies in one session, so the unit is a matched pair and the count follows McNemar, which depends
on the discordance rate. The share of initial conditions on which the two policies disagree, and
not on the two rates alone. Write $\delta = |p_{1} - p_{2}|$ for the gap between the two success
rates, $\pi_{d}$ for that discordance rate, $z_{1-\alpha/2} = 1.96$ at $\alpha = 0.05$ and
$z_{1-\beta} = 0.84$ at 80 percent power; then the number of matched pairs is
\begin{equation}
n_{\mathrm{pairs}} =
  \frac{\bigl[\,z_{1-\alpha/2}\sqrt{\pi_{d}}
              + z_{1-\beta}\sqrt{\pi_{d} - \delta^{2}}\,\bigr]^{2}}{\delta^{2}}.
\label{eq:mcnemar}
\end{equation}
To separate 50 from 70 percent, which is $\delta = 0.20$, at $\alpha = 0.05$ with 80 percent power,
\eqref{eq:mcnemar} asks per arm for
\begin{equation*}
\begin{array}{lrrrr}
\toprule
\text{discordance } \pi_{d} & 0.2 & 0.3 & 0.4 & 0.5\\
\midrule
n_{\mathrm{pairs}} & 37 & 57 & 77 & 96\\
\bottomrule
\end{array}
\end{equation*}
The protocol assumes 0.3 and asks for 57, and states the sensitivity rather than hiding it, because
0.5 is the discordance the same two rates produce when the pairing buys nothing, and at that value
the paired count returns to the 93 per arm an unpaired test would need. The saving from pairing is
real but smaller than the pair counts suggest, since a pair costs two rollouts: 57 pairs is 114
rollouts against 186. An earlier version of this section quoted 93, 169 and 387 per arm for gaps of
20, 15 and 10 points, which are correct for independent arms and are the wrong test for this
protocol. That 93 was also the same integer as the half-width calculation in the paragraph above,
which is a coincidence of the worst-case arithmetic and not a second derivation of the same number.

One hundred is a cap and not a bill, because on a graded score a sequential test reached its
decision in 12 to 36 paired hardware trials in \cite{beyond_binary_success_2026}. At 30 rollouts a
continuous score already carries a half-width of $\pm 0.36$ standard deviations, which is why a
graded score can stop where a binary one cannot. A cell that stops early does not report a Wilson
interval. Optional stopping breaks the coverage of a fixed-$n$ interval, which is the reason
Snyder et al. \cite{beyond_binary_success_2026} and Badithela et al. \cite{suresim_2025} use
anytime-valid betting intervals rather than Wilson, so Table~\ref{tab:protocol} asks a cell run to
a fixed 100 for a Wilson interval and a cell stopped early for a confidence sequence, and never
for both. Intervals are also marginal rather than simultaneous. For example, evaluating twelve methods on seven axes produces 84 intervals. At 84 independent 95 percent intervals, four
excursions are expected by construction, so a paper comparing $k$ policies on $m$ tasks corrects
its $k(k-1)/2$ pairwise tests to a global 95 percent level, as the TRI LBM Team
\cite{lbm_careful_examination_2025} does, or says its intervals are not simultaneous.

Simulation is cheap, so simulated cells take 200 episodes, giving a 6.9-point half-width.
Perturbation axes are screened rather than certified, and 40 per axis buys a 15-point half-width
on each axis's own absolute rate, which is enough to rank the axes and pick the two worst for
hardware. It is not enough for the ratio to the anchor that an earlier draft asked each cell to
report. At 40 trials in each arm, a fall from a 0.50 anchor to 0.30 is a ratio of 0.60 with a 95
percent interval of 0.34 to 1.06, which contains 1: at the screening count you cannot establish
that the perturbation hurt at all. Certifying that same drop takes 101 per arm at 80 percent
power, or about 50 for an interval that merely excludes 1, so Table~\ref{tab:protocol} now asks for
absolute rates with their own intervals at 40, reports the ratio without an interval, and prescribes
101 before any claim that a named axis hurt.

For unseen objects the resampling unit is the object and not the trial, so 20 objects at 5 trials
each gives 100 trials and an object-level half-width near 20 points. That 20 points is the Wilson
width at $n = 20$ and it treats each object's outcome as a single Bernoulli draw, which the five
within-object trials are not. It is the right order of magnitude and the assumption belongs in the
cell. A 10-point claim about an object distribution needs about 93 objects. Seven of the 39 rows
that state an unseen count reach that: 225 in \cite{bimangrasp_2024}, 241 in \cite{resdex_2024}
and \cite{unidexgrasp_2023}, 360 in \cite{dexgraspvla_2025}, 500 in \cite{dexmv_2021}, 2029 in
\cite{clutterdexgrasp_2025} and 503409 in \cite{graspxl_2024}. For a continuous score the
half-width is 1.96 standard deviations over the square root of the count, $1.96\,s/\sqrt{n}$, so
100 rollouts give $\pm 0.20$ standard deviations, and the unit is the rollout because frames
within one are correlated.

The transfer axis is the one where 100 is least defensible. On 100 matched pairs a measured
correlation of 0.70 carries a Fisher-z interval of 0.58 to 0.79, a half-width of about 0.10 that
130 pairs would be needed to guarantee. That is enough to establish that a simulator tracks
reality at all, and it is not enough to separate \cite{suresim_2025}'s useful regime from its
marginal one, since those differ by about 0.11 in correlation. Both limits come within 0.05 of the
estimate only at about 457 pairs. Table~\ref{tab:protocol} states which of the two decisions each
count supports rather than leaving a reader to assume the larger one.
